% !TeX program = lualatex
% !TeX encoding = utf8
% !TeX spellcheck = uk_UA
% !TeX root =../PyscfBook.tex


%% --------------------------------------------------------
\section{$T_1$-діагностика}
%% --------------------------------------------------------

Однореферентні методи, такі як \texttt{MP2}, \texttt{CCSD} чи \texttt{CCSD(T)},
передбачають, що хвильова функція системи добре описується
одним детермінантом Гартрі–Фока. Проте у випадках, коли
у системі присутня \textbf{статична (нединамічна) кореляція} ---
наприклад, при розриві зв’язку, у відкритих оболонках або при
наявності кількох близьких за енергією конфігурацій ---
таке наближення стає непридатним.

Для кількісного оцінювання ступеня мультиреферентності
використовується \textbf{$T_1$-діагностика} (\emph{$T_1$ diagnostic}),
запропонована Лі та Таулером (Lee \& Taylor, 1989).
Вона визначається через амплітуди одиночних збуджень \( t_i^a \)
у методі \texttt{CCSD}:

\begin{equation}
    T_1 = \sqrt{\frac{1}{n_\text{occ}} \sum_{i,a} (t_i^a)^2},
\end{equation}
де сума ведеться за всіма зайнятими (\(i\)) та віртуальними (\(a\))
орбіталями, а \(n_\text{occ}\) --- число зайнятих орбіталей.

\paragraph{Фізичний зміст.}
Амплітуди \(t_i^a\) характеризують внесок одноелектронних
збуджень у хвильову функцію \(|\Psi_{\text{CC}}\rangle = e^{\hat{T}}|\Phi_0\rangle\).
Якщо система справді однореферентна, ці амплітуди дуже малі,
оскільки основна конфігурація добре описує стан системи.
Навпаки, великі \(t_i^a\) свідчать про значне змішування детермінантів
і, відповідно, про порушення однодетермінантного опису.

\paragraph{Типові орієнтири.}
Для закритооболонкових молекул, які адекватно описуються
методами \texttt{MP2} або \texttt{CCSD}, характерні значення
\[
T_1 < 0.02.
\]
Якщо ж \(T_1 > 0.02\), система, ймовірно, має мультиреферентний характер,
і результати однодетермінантних методів можуть бути ненадійними.
У таких випадках рекомендується застосовувати багатореферентні методи,
зокрема \texttt{CASSCF} або \texttt{NEVPT2}.

\paragraph{Практичне обчислення.}
У пакеті \texttt{PySCF} значення \(T_1\) можна отримати
безпосередньо після виконання розрахунку \texttt{CCSD}:
\begin{verbatim}
from pyscf import cc
mycc = cc.CCSD(mf).run()
T1 = mycc.t1_diagnostic()
print("T1 diagnostic =", T1)
\end{verbatim}

\paragraph{Приклад.}
Для молекули \ce{HF} у базисі \texttt{cc-pVDZ} значення \(T_1 \approx 0.011\),
що свідчить про відсутність суттєвої статичної кореляції
і підтверджує придатність методів \texttt{CCSD} та \texttt{CCSD(T)}.

\begin{remark}
$T_1$-діагностика --- це не формальний критерій збіжності,
а емпіричний показник, який допомагає оцінити
адекватність однодетермінантного підходу.
Для більш складних випадків також застосовують
індикатор $D_1$, побудований на базі власних значень
матриці \( \mathbf{t}_1 \mathbf{t}_1^\mathrm{T} \).
\end{remark}
